\clearpage
\fancyhead[L]{\small VESTA.}
\hypertarget{native-295-lower}{}
\centerline{VESTA.}

\begin{center}
\rule{0.5\linewidth}{0.4pt}\\[2pt]
Monatliche Correspondenz. Band XVI. S. 285–291. September 1807.\\[2pt]
\rule{0.5\linewidth}{0.4pt}
\end{center}

\smallskip
\noindent---\ ---\ ---

\smallskip
\par\indent „Am 8. Juli, schreibt uns Dr. \textsc{Olbers}, habe ich meine Beobachtungen abgeschlossen, weil der tiefe Stand des Planeten, der Mondschein und die Dämmerung sie zu sehr erschwerten. Die letzten sechs Beobachtungen wurden in Gegenwart unseres vortrefflichen Freundes \textsc{Gauss} gemacht, der mich mit einem Besuch erfreut hat. Dieser ausgezeichnete Astronom wird uns nun sehr bald die III. Elemente der \emph{Vesta} geben. Die zweiten Elemente werden unterdessen nur geringfügige Korrekturen erfahren, da sie auch die letzten Beobachtungen noch so genau wiedergeben. Die Rektaszension ist nach ihnen nur etwa $1'$ zu groß und die Deklination um weniger als $30''$ zu klein. Aus den II. Elementen habe ich ermittelt, dass \emph{Vesta} am 25. März 1796 in der Zone stand, die laut der \emph{Hist. Cél.} damals beobachtet wurde. Doch kann ein kleiner, durchaus
\SourceLeaf{296}möglicher Fehler in der großen Halbachse und damit in der mittleren Bewegung der \emph{Vesta} ihre Deklination für diesen Tag noch stark verändern. Ob jetzt ein beobachteter Stern in dieser Zone fehlt, weiß ich nicht, denn der Teil, in dem die \emph{Vesta} gestanden haben könnte, ist jetzt im Licht der Sonne verborgen.“

\par\indent Wir haben diese \emph{dritten Elemente} vor wenigen Tagen erhalten und geben sie hier wieder. Dr. \textsc{Gauss} hatte dafür außer allen in dieser Zeitschrift schon veröffentlichten Beobachtungen auch eine Reihe Pariser Beobachtungen vom 13. April bis zum 22. Mai verwendet, die ihm Herr \textsc{Bouvard} mitgeteilt hatte.

\smallskip

\begin{center}\begin{tabular}{lr}
Epoche 31. März 1807, Mittag in Bremen  & $192^\circ\;23'\;30''1$ \\
Tägliche tropische Bewegung .\ .\ . & $981''8459$ \\
\shortstack[l]{Aphel, für die Epoche siderisch\\fest} .\ .\ .\ .\ .\ .\ .\ .\ .\ . & $69^\circ\;50'\;31''9$ \\
Aufsteigender Knoten .\ .\ .\ .\ . & $103^\circ\;18'\;28''$ \\
Neigung der Bahn .\ .\ .\ .\ .\ .\ . & $7\;8\;10.7$ \\
Exzentrizität .\ .\ .\ .\ .\ .\ .\ .\ .\ . & $0.0855050$ \\
Logarithmus des mittleren Abstands .\ . & $0.3720160$ \\
\end{tabular}\end{center}
\par\indent Diese Elemente unterscheiden sich nur geringfügig von den zweiten. Dadurch wird jedoch die geozentrische Position der \emph{Vesta} zur Zeit der nächsten Opposition im September 1808 um $22'$ größer. Zugleich hatte Dr. \textsc{Gauss} die Güte, uns jene Reihe Pariser Beobachtungen samt den Differenzen zwischen den beobachteten und den nach den II. Elementen berechneten Positionen zu schicken.

\smallskip
\clearpage\begingroup\fontsize{9}{11}\selectfont\setlength{\tabcolsep}{2.5pt}\renewcommand{\arraystretch}{1.05}
\begin{center}\begin{tabular}{lrrrrr}
1807 & \shortstack{Mittlere Zeit\\in Paris} & Rektasz. \VestaRATwoNinetySix & \shortstack{Differenz zu\\II. Elem.} & Dekl. \VestaDecTwoNinetySix & \shortstack{Differenz zu\\II. Elem.} \\ \hline
April 13 & $22^{\text{u}}\;39^{\text{m}}\;6\mathord{\raisebox{.65ex}{\ObservedTimeStroke}}2$ & $180^\circ 57' 7''8$ & $+3''4$ & $12^\circ 50' 35''0$ & $+5''2$ \\
14 & $22\;34\;28.3$ & $180\;46\;27.9$ & $+1.3$ & $12\;52\;39.2$ & $+6.0$ \\
18 & $22\;16\;7.0$ & $180\;\;7\;\;6.9$ & $-3.1$ & $12\;58\;\;9.3$ & $+12.7$ \\
19 & $22\;11\;34.9$ & $179\;58\;\;4.1$ & $+2.2$ & $12\;59\;\;5.5$ & $+0.4$ \\
21 & $22\;2\;39.0$ & $179\;41\;20.6$ & $-1.4$ & $12\;58\;40.7$ & $+65.0$ \\
24 & $21\;49\;20\;2$ & $179\;19\;\;5.1$ & $+1.7$ & $12\;58\;44.2$ & $+1.7$ \\
25 & $21\;44\;58.3$ & $179\;12\;26.6$ & $+4.7$ & $12\;57\;47.3$ & $+7.4$ \\
26 & $21\;40\;38.5$ & $179\;\;6\;33.8$ & $-13.1$ & $12\;56\;40.6$ & $+7.6$ \\
27 & $21\;36\;18.4$ & $179\;\;0\;35.4$ & $+0.1$ & $12\;55\;10.8$ & $+15.9$ \\
28 & $21\;32\;1.2$ & $178\;55\;15.2$ & $-0.$ & $12\;53\;47.3$ & $+2.8$ \\
29 & $21\;27\;45.5$ & $178\;50\;16.5$ & $+4.1$ & $12\;51\;54.4$ & $+4.3$ \\
30 & $21\;23\;31.7$ & $178\;45\;52.5$ & $-0.8$ & $12\;49\;49.4$ & $+3.2$ \\
Mai 1 & $21\;19\;20.0$ & $178\;41\;47.7$ & $+1.6$ & $12\;47\;26.9$ & $+5.0$ \\
7 & $20\;54\;44.7$ & $178\;26\;43.7$ & $+4.3$ & $12\;28\;39.6$ & $-7.5$ \\
12 & $20\;35\;3.3$ & $178\;26\;18.3$ & $+8.1$ & $12\;\;6\;35.8$ & $+2.8$ \\
15 & $20\;23\;35.2$ & $178\;31\;18.5$ & $+7.9$ & $11\;50\;39.4$ & $+25.5$ \\
16 & $20\;19\;50.7$ & $178\;33\;47.4$ & $+9.8$ & $11\;45\;30.8$ & $+0.2$ \\
17 & $20\;16\;5.0$ & $178\;36\;42.3$ & $+10.9$ & $11\;39\;46.8$ & $-0.6$ \\
18 & $20\;13\;22.5$ & $178\;39\;57.3$ & $+16.6$ & $11\;33\;50.4$ & $+0.3$ \\
20 & $20\;5\;1.9$ & $178\;47\;55.5$ & $+13.0$ & $11\;21\;31.8$ & $-3.5$ \\
21 & $20\;1\;24.6$ & $178\;52\;29.4$ & $+12.4$ & $11\;15\;\;5.5$ & $-3.6$ \\
22 & $19\;57\;48.9$ & $178\;57\;25.8$ & $+12.9$ & $11\;\;8\;27.6$ & $-2.1$ \\
\end{tabular}\end{center}\endgroup


\par\indent Die Zeiten werden hier von Mitternacht an gezählt, und die Deklination vom 21. April ist offenbar um eine Bogenminute zu klein. Doch wer staunt nicht über diese ganz ausgezeichnete Übereinstimmung der II. Elemente mit dem Himmel?

\par\indent „Die Ephemeride für den Lauf der \emph{Vesta} im nächsten Jahr, schreibt uns Dr. \textsc{Gauss} weiter, lege ich bei. Ein großer Teil des diesmaligen Laufs der \emph{Vesta} fällt auf die Karte von \textsc{Harding} für die Bewegung der \emph{Ceres} im Jahr 1804, die der M. Corr. beigefügt ist. Professor \textsc{Harding} wird das Blatt seines Atlas, das den übrigen Teil enthält, demnächst dem Kupferstecher übergeben. Dieses nützliche Vorhaben kann man gar nicht genug empfehlen und unterstützen. Das erste, sehr sauber gestochene Blatt habe ich bereits in \emph{Bremen} gesehen.“

\SourceLeaf{297}\par\indent Leider wird die \emph{Vesta} nach der beiliegenden Ephemeride erst im Mai oder Juni des folgenden Jahres wieder sichtbar. [Lauf der Vesta vom 2. Mai 1808 bis zum 18. März 1809 in Abständen von vier Tagen nach den III. Elementen. Rektaszension und Deklination auf Bogenminuten, Logarithmus des Abstands auf vier Dezimalstellen.]

\par\indent Da die Elemente der \emph{Vesta}-Bahn jetzt schon recht gut angenähert sind, wird es interessant sein, in älteren Sternverzeichnissen nach \emph{Vesta} zu suchen. Unser Versuch, eine Beobachtung der \emph{Vesta} unter den Beobachtungen von \textsc{Mayer} zu finden, war bisher vergeblich.

\begin{flushright}
[\textsc{von Zach}.]
\end{flushright}

